\section{1次分数変換}\label{節：一次分数変換}
\subsection{不動点を引く理由}

次の項が
\[
 a_{n+1}=f(a_n),\qquad
 f(z)=\frac{pz+q}{rz+s}
\]
によって決まる場合を考える.
この形の$f$を\textbf{1次分数変換}という.
本節では係数は実数,\,$r\ne0$,\,$ps-qr\ne0$とし,\,元の更新で分母が零にならない範囲を扱う.
$r=0$なら1階の線形の式へ戻り,\,$ps-qr=0$なら定義できる入力で右辺が定数になるため,\,別に調べればよい.

不動点$\alpha$は
\[
 \alpha=\frac{p\alpha+q}{r\alpha+s}
\]
を満たすので,
\[
 r\alpha^2+(s-p)\alpha-q=0
\]
から求まる.
この方程式の根では$r\alpha+s\ne0$である.
もし零なら,\,元の方程式から$p\alpha+q=0$にもなり,\,$ps-qr\ne0$に反する.

不動点の式を使って引き算すると
\[
 f(z)-\alpha
 =\frac{(p-r\alpha)(z-\alpha)}{rz+s}.
\]
定数項は消えるが,\,まだ分母に$z$が残る.
ここで,\,異なるもう一つの不動点$\beta$があれば
\[
 f(z)-\beta
 =\frac{(p-r\beta)(z-\beta)}{rz+s}.
\]
二式の比を取れば,\,共通の分母が消える.
\[
 \frac{f(z)-\alpha}{f(z)-\beta}
 =\frac{p-r\alpha}{p-r\beta}
  \frac{z-\alpha}{z-\beta}.
\]
つまり
\[
 b_n=\frac{a_n-\alpha}{a_n-\beta},\qquad
 \kappa=\frac{p-r\alpha}{p-r\beta}
\]
と置くと,\,$b_{n+1}=\kappa b_n$である.
ここで$p-r\alpha,\,p-r\beta$も非零である.
実際,\,不動点の式から
$(p-r\alpha)(r\alpha+s)=ps-qr$となる.

初項が不動点なら定数列として先に扱う.
$a_1\ne\beta$なら,\,上の差の式から$a_n\ne\beta$が続くため,\,比を取れる.
解いた後は
\[
 b_n=b_1\kappa^{n-1},\qquad
 a_n=\frac{\beta b_n-\alpha}{b_n-1}
\]
で元へ戻す.
$b_n=1$になる場合は有限の$a_n$へ戻せず,\,元の更新もその項まで実数列としては定義できない.

\begin{bookpoint}
一つの不動点を引くと定数項が消える.
二つの不動点からの差の比を取ると,\,残った共通の分母も消える.
置き換えは,\,消したい部分に合わせて組み立てられる.
\end{bookpoint}

\noindent\textbf{【例題】}\par\nobreak
$a_1=3$,\,$a_{n+1}=\frac{2a_n+1}{a_n+2}$を解け.

\noindent\textbf{【方針】}\par\nobreak
不動点1と$-1$からの差の比を取ります.

\noindent\textbf{【解答】}\par\nobreak
\[
 a_n=\frac{2\cdot3^{n-1}+1}{2\cdot3^{n-1}-1}.
\]
\noindent\textbf{【解法】}\par\nobreak
$b_n=\frac{a_n-1}{a_n+1}$とすると
$b_{n+1}=\frac13b_n$,\,$b_1=\frac12$である.
従って$b_n=\frac1{2\cdot3^{n-1}}$となり,
$a_n=\frac{1+b_n}{1-b_n}$から答えを得る.
分母は正で,\,各項も正なので元の分母$a_n+2$は零にならない.
また$b_n\to0$だから$a_n\to1$となる.
不動点からの距離そのものではなく,\,二つの差の比が縮小していた.

\subsection{不動点が重なる場合}

不動点の方程式が重解$\alpha$を持つと,\,異なる二つの差を作れない.
この場合は
\[
 2r\alpha=p-s,\qquad
 p-r\alpha=r\alpha+s=\ell
\]
となり,\,$\ell\ne0$である.
$u_n=a_n-\alpha$と置くと
\[
 u_{n+1}=\frac{\ell u_n}{ru_n+\ell}.
\]
初項が不動点の場合を除き,\,逆数を取って
\[
 \frac1{u_{n+1}}=\frac1{u_n}+\frac r\ell
\]
という等差型へ直せる.
二つの固有方向が分かれない場合に足し算が残った第2章と,\,似た現象が現れた.
その直接の対応は最後の節で確認する.

例えば
\[
 a_1=5,\qquad a_{n+1}=\frac{4a_n-9}{a_n-2}
\]
では$\alpha=3$であり,
\[
 \frac1{a_{n+1}-3}=\frac1{a_n-3}+1.
\]
従って$a_n=3+\frac2{2n-1}$となる.
この例では$a_n>3$が保たれ,\,元の分母も非零である.

\subsection{実数の数列を単位円へ}

不動点の二解が非実数なら,\,実係数なので$\alpha,\overline\alpha$という共役の組になる.
実数$a_n$について
\[
 b_n=\frac{a_n-\alpha}{a_n-\overline\alpha}
\]
と置く.
分子と分母は共役だから,\,$|b_n|=1$である.
また
\[
 \kappa=\frac{p-r\alpha}{p-r\overline\alpha}
\]
も分子と分母が共役なので$|\kappa|=1$となる.
従って,\,$b_{n+1}=\kappa b_n$は\textbf{単位円周上の一定角の回転}である.
実軸上で複雑に見えた動きが,\,円周上では単純になる.

\noindent\textbf{【例題】}\par\nobreak
\[
 a_1=\frac12,\qquad a_{n+1}=\frac{a_n+1}{1-a_n}
\]
を複素平面で解釈し,\,周期を求めよ.

\noindent\textbf{【方針】}\par\nobreak
不動点は$i,-i$です. 差の比を取り,\,一回の回転角を調べます.

\noindent\textbf{【解答】}\par\nobreak
最小周期は4であり,
\[
 \frac12,\ 3,\ -2,\ -\frac13,\ \frac12,\ldots
\]
と繰り返す.

\noindent\textbf{【解法】}\par\nobreak
$b_n=\frac{a_n-i}{a_n+i}$とすると
\[
 b_{n+1}=\frac{1+i}{1-i}b_n=ib_n.
\]
毎回反時計回りに$\frac{\pi}{2}$回転する.
$b_n$は零でなく,\,4回より少ない正の回数では元へ戻らない.
復元は$a_n=i\frac{b_n+1}{1-b_n}$である.
実際に元の式で四項を計算すると上の並びを得る.
どの項も1ではないため,\,すべての更新が定義できる.

ここで,\,回転が定義できることと,\,元の実数列が全項で定義できることは区別する.
例えば同じ式で$a_1=0$とすると$a_2=1$となり,\,$a_3$は定義できない.
円周上では回転を続けられるが,\,復元式の分母が零になる点に到達している.

\subsection{実軸と円を結ぶ座標}\label{節：実軸と円}

回転を三角関数と結び付けるため,\,別の一定倍をした座標
\[
 w=\Phi(a)=\frac{1+ia}{1-ia}
\]
を使おう.
実数$a$なら$|w|=1$であり,\,この$\Phi(a)$は先ほどの$b$に対して$w=-b$である.
例えば$a=0$は$w=1$,\,$a=1$は$w=i$へ移る.

$a=\tan\theta$と書ける範囲では,
\[
 \begin{aligned}
 \Phi(\tan\theta)
 &=\frac{\cos\theta+i\sin\theta}{\cos\theta-i\sin\theta}\\
 &=\cos2\theta+i\sin2\theta.
 \end{aligned}
\]
実数$a$を傾きとして記録した角度が,\,円周上ではその2倍の角度になる.
\NewFigCayley

$w=u+iv$が単位円上で$w\ne-1$なら,\,逆に
\[
 a=\frac{v}{1+u}
\]
で戻れる.
点$(-1,0)$から$(u,v)$への直線の傾きが$a$である.
従って$\Phi$は実軸と,\,一点$-1$を除いた単位円を一対一に対応させる.
除いた一点は,\,傾きが無限に大きくなる方向に当たる.
必要に応じて実軸に無限遠点$\infty$を一つ加えると円全体と対応するが,\,実数の割り算では依然として分母の条件が必要である.

実数$t$を固定し,
\[
 f(a)=\frac{a+t}{1-ta}
\]
とすると,\,直接の計算から
\[
 \Phi(f(a))=\Phi(t)\Phi(a)
\]
となる. $t=\tan\phi$なら$\Phi(t)=e^{2i\phi}$であり,
元の式は$a=\tan\theta$の角度を$\theta+\phi$へ進める式である.
正接の加法公式
\[
 \tan(\theta+\phi)
 =\frac{\tan\theta+\tan\phi}{1-\tan\theta\tan\phi}
\]
が,\,円周上の掛け算として現れた.
角度は$\pi$の整数倍だけ違っても同じ傾きを表す.
